% ============================================================
% 30_layout_spacing.tex -- Layout-Grundlagen und Spacing-Skala
% ============================================================

\setlength{\parskip}{1.2pt plus 0.8pt minus 0.2pt}
\setlength{\parindent}{0pt}

\widowpenalty=10000
\clubpenalty=10000
\brokenpenalty=4000
\predisplaypenalty=10000

\raggedcolumns
\newcounter{zsfFirstChap}\setcounter{zsfFirstChap}{1}

% Kompakter als Template, ruhiger als die MaschKo-Maximalverdichtung.
\newlength{\ZSFspaceXXS}\setlength{\ZSFspaceXXS}{0.5pt}
\newlength{\ZSFspaceXS} \setlength{\ZSFspaceXS}{1pt}
\newlength{\ZSFspaceS}  \setlength{\ZSFspaceS}{2.5pt}
\newlength{\ZSFspaceM}  \setlength{\ZSFspaceM}{5pt}
\newlength{\ZSFspaceL}  \setlength{\ZSFspaceL}{7pt}

\newlength{\ZSFcornerBox}   \setlength{\ZSFcornerBox}{0.7mm}
\newlength{\ZSFcornerInline}\setlength{\ZSFcornerInline}{1mm}

\setlength{\columnsep}{\ZSFspaceL}
\setlength{\multicolsep}{\ZSFspaceS}
\setlength{\emergencystretch}{2.5em}

% Vertikaler Rhythmus: ein einziges einfaches System.
% Bars und Boxen bringen ihre eigenen Skips mit. Titelbalken reservieren
% nur genug Folgeraum, damit sie nicht allein am Spaltenende stehen.
\newlength{\ZSFtitleBarHeight}\setlength{\ZSFtitleBarHeight}{7.2pt}
\newlength{\ZSFchapterBeforeSkip}\setlength{\ZSFchapterBeforeSkip}{\ZSFspaceM}
\newlength{\ZSFchapterAfterSkip}\setlength{\ZSFchapterAfterSkip}{\ZSFspaceXS}
\newlength{\ZSFsubsectionBeforeSkip}\setlength{\ZSFsubsectionBeforeSkip}{\ZSFspaceS}
\newlength{\ZSFsubsectionAfterSkip}\setlength{\ZSFsubsectionAfterSkip}{\ZSFspaceXS}
\newlength{\ZSFsubsubsectionBeforeSkip}\setlength{\ZSFsubsubsectionBeforeSkip}{\ZSFspaceXS}
\newlength{\ZSFsubsubsectionAfterSkip}\setlength{\ZSFsubsubsectionAfterSkip}{0pt}
\newlength{\ZSFboxBeforeSkip}\setlength{\ZSFboxBeforeSkip}{\ZSFspaceXS}
\newlength{\ZSFboxAfterSkip}\setlength{\ZSFboxAfterSkip}{\ZSFspaceS}
\newlength{\ZSFbreakThreshold}\setlength{\ZSFbreakThreshold}{3\baselineskip}

% Display-Mathematik folgt derselben kompakten Rhythmik wie Boxen und Balken.
% Die Zuweisung im Begin-Document-Hook gewinnt gegen die Klassen-Defaults.
\newlength{\ZSFdisplayBeforeSkip}     \setlength{\ZSFdisplayBeforeSkip}{\ZSFspaceS}
\newlength{\ZSFdisplayAfterSkip}      \setlength{\ZSFdisplayAfterSkip}{\ZSFspaceS}
\newlength{\ZSFdisplayShortBeforeSkip}\setlength{\ZSFdisplayShortBeforeSkip}{\ZSFspaceXS}
\newlength{\ZSFdisplayShortAfterSkip} \setlength{\ZSFdisplayShortAfterSkip}{\ZSFspaceXS}
\newlength{\ZSFmathLineGap}           \setlength{\ZSFmathLineGap}{\ZSFspaceXS}

\newcommand{\ZSFApplyDisplaySpacing}{%
  \setlength{\abovedisplayskip}{\ZSFdisplayBeforeSkip}%
  \setlength{\belowdisplayskip}{\ZSFdisplayAfterSkip}%
  \setlength{\abovedisplayshortskip}{\ZSFdisplayShortBeforeSkip}%
  \setlength{\belowdisplayshortskip}{\ZSFdisplayShortAfterSkip}%
  \setlength{\jot}{\ZSFmathLineGap}%
}
\AtBeginDocument{\ZSFApplyDisplaySpacing}

\newcommand{\ZSFHeadingBefore}[1]{\par\needspace{\ZSFbreakThreshold}\addvspace{#1}}
\newcommand{\ZSFHeadingAfter}[1]{\par\nobreak\vspace{#1}\nobreak}

\newlength{\ZSFborderThin}  \setlength{\ZSFborderThin}{0.35pt}
\newlength{\ZSFborderBox}   \setlength{\ZSFborderBox}{0.45pt}
\newlength{\ZSFborderStrong}\setlength{\ZSFborderStrong}{0.8pt}

\newlength{\ZSFstatementBarWidth}\setlength{\ZSFstatementBarWidth}{2pt}
\newlength{\ZSFstatementLeftPad} \setlength{\ZSFstatementLeftPad}{\ZSFspaceM}
\newlength{\ZSFprocStepGap}      \setlength{\ZSFprocStepGap}{\ZSFspaceXS}
\newlength{\ZSFprocItemSep}      \setlength{\ZSFprocItemSep}{\ZSFspaceXS}
\newlength{\ZSFprocLeftMargin}   \setlength{\ZSFprocLeftMargin}{1.6em}
\newlength{\ZSFlistLeftMargin}   \setlength{\ZSFlistLeftMargin}{1.2em}
\newlength{\ZSFprocLabSep}       \setlength{\ZSFprocLabSep}{0.4em}

\newlength{\ZSFcodeMiniSkip} \setlength{\ZSFcodeMiniSkip}{0.4mm}
\newlength{\ZSFcodeFrameRule}\setlength{\ZSFcodeFrameRule}{\ZSFborderThin}
\newlength{\ZSFcodeInnerXPad}\setlength{\ZSFcodeInnerXPad}{\ZSFspaceS}
\newlength{\ZSFcodeInnerYPad}\setlength{\ZSFcodeInnerYPad}{\ZSFspaceXXS}

% Dimensionslose Tabellen-Rhythmik; ausgewählt über die ZSFtable*-API.
\newcommand{\ZSFtableRowStretchDefault}{1}
\newcommand{\ZSFtableRowStretchRoomy}{1.25}

% Standard-Breite für Grafiken (z.B. 2/3 der Spaltenbreite)
\newlength{\ZSFgraphicswidth}
\setlength{\ZSFgraphicswidth}{0.667\columnwidth}


\newlength{\ZSFcodeCommentMaxWidthDefault}
  \setlength{\ZSFcodeCommentMaxWidthDefault}{.25\columnwidth}
\newlength{\ZSFcodeCommentThresholdDefault}
  \setlength{\ZSFcodeCommentThresholdDefault}{\maxdimen}
\newlength{\ZSFcodeCommentMaxWidth} \setlength{\ZSFcodeCommentMaxWidth}{\ZSFcodeCommentMaxWidthDefault}
\newlength{\ZSFcodeCommentThreshold}\setlength{\ZSFcodeCommentThreshold}{\ZSFcodeCommentThresholdDefault}
\newlength{\ZSFcodeCommentGap}      \setlength{\ZSFcodeCommentGap}{\ZSFspaceM}

\setlist{nosep, itemsep=0pt, parsep=0pt, topsep=\ZSFspaceXS, partopsep=0pt}
\setlist[itemize]{leftmargin=3mm}
\setlist[enumerate]{leftmargin=4mm}

\newcommand{\ZSFgapXS}{\vspace{\ZSFspaceXS}}
\newcommand{\ZSFgapS}{\vspace{\ZSFspaceS}}
\newcommand{\ZSFgapM}{\vspace{\ZSFspaceM}}
\newcommand{\ZSFgapL}{\vspace{\ZSFspaceL}}
\newcommand{\ZSFSectionGap}{\vspace{\ZSFspaceL}}
\newcommand{\ZSFtabColGap}{\hspace{1.5em}}
\newcommand{\ZSFmanualColumnBreak}{\columnbreak}
